% Generated draft: requires mathematical and source review.
\documentclass[11pt]{article}
\usepackage[margin=25mm]{geometry}
\usepackage{fontspec}
\setmainfont{Latin Modern Roman}
\usepackage{amsmath,amssymb,mathtools,unicode-math}
\setmathfont{Latin Modern Math}
\usepackage{enumitem,needspace,xurl,hyperref}
\setlength{\parindent}{0pt}
\setlength{\parskip}{5pt}
\setlength{\emergencystretch}{3em}
\newcommand{\N}{\mathbb{N}}
\newcommand{\Z}{\mathbb{Z}}
\newcommand{\Q}{\mathbb{Q}}
\newcommand{\R}{\mathbb{R}}
\newcommand{\C}{\mathbb{C}}
\newcommand{\F}{\mathbb{F}}
\newcommand{\A}{\mathbb{A}}
\newcommand{\PP}{\mathbb{P}}
\newcommand{\E}{\mathbb{E}}
\newcommand{\eps}{\varepsilon}
\newcommand{\dd}{\,\mathrm{d}}
\begin{document}
\section*{UnsolvedMath ID: 800961}
\textbf{Source problem code:} catalogue entry
\subsection{Definitions and mathematical statement}
Let \(M\) be a smooth manifold of dimension \(n\). A Finsler metric on \(M\) is a continuous function \(F:TM\to[0,\infty)\), smooth away from the zero section, positive on nonzero tangent vectors, positively homogeneous of degree one, and strongly convex: for each \(y\in T_xM\setminus\{0\}\), the fibre Hessian
\[
g_y=\tfrac12 D_y^2(F^2)
\]
is positive definite. Reversibility, meaning \(F(x,y)=F(x,-y)\), is not included in this definition. Positive flag curvature means
\[
K_F(x,y,P)>0
\]
for every \(x\in M\), every two-dimensional subspace \(P\subseteq T_xM\), and every nonzero flagpole \(y\in P\).

A manifold is closed if it is compact and has no boundary. For an embedded smooth submanifold \(N\subseteq M\), total geodesy means that every geodesic of the induced Finsler metric \(F|_{TN}\) is also an ambient geodesic. Equivalently, an ambient geodesic initially tangent to \(N\) remains in \(N\) locally.

\textbf{Question (Burns--Matveev, Question 42).} Must any two nonempty closed totally geodesic submanifolds \(N_1,N_2\) of a closed Finsler manifold \((M,F)\) with positive flag curvature satisfy
\[
\dim N_1+\dim N_2\geq n
\quad\Longrightarrow\quad
N_1\cap N_2\neq\varnothing?
\]
This asks whether the indicated conclusion of Frankel's theorem extends to Finsler manifolds.

\textbf{Unresolved conventions.} The supplied statement does not explicitly specify connectedness of the ambient manifold or the submanifolds, embedded versus immersed submanifolds, or reversibility of the metric. The displayed formulation uses the conventional nonempty embedded interpretation of submanifolds and the general, possibly nonreversible, definition of a Finsler metric. These interpretations require confirmation against the source. No connectedness hypothesis has been added; whether connectedness is implicit in the intended source formulation remains unresolved.
\subsection{Short English statement}
Burns and Matveev ask whether two closed totally geodesic submanifolds in a closed positively flag-curved Finsler manifold must intersect whenever their dimensions sum to at least the ambient dimension. The supplied registry labels the question open; this status is not independently verified. Connectedness, reversibility, and submanifold conventions require review.
\subsection{Sources}
\begin{itemize}
\item Local UnsolvedMath registry: permanent ID \(800961\), source problem code \texttt{catalogue entry}, title \emph{Open Problems and Questions about Geodesics --- Question 42}. The supplied record attributes the question to Burns and Matveev, records proposed year 2019, and labels its status \texttt{open}.
\item Burns and Matveev, \emph{Open Problems and Questions about Geodesics} (2019, as identified in the supplied source list), Question 42. Supplied source URL: \url{https://arxiv.org/abs/1308.5417}. The provenance records access on 2026-10-16 and extraction by \texttt{arxiv\_latex}; these are supplied provenance claims, not an assertion of independent source access or verification.
\end{itemize}
\end{document}
